\PassOptionsToPackage{dvipsnames,svgnames,table}{xcolor}
\documentclass[10pt]{article}
\usepackage[a4paper,margin=22mm]{geometry}
\usepackage[no-math]{fontspec}\setmainfont[SmallCapsFont={Latin Modern Roman Caps}]{Latin Modern Roman}
\usepackage{amsmath,amssymb,amsthm,mathtools,graphicx,xcolor,hyperref,xurl,xspace,ifthen,enumerate,textcomp}
\usepackage[shortlabels]{enumitem}
\setlength{\emergencystretch}{4em}
\newtheorem{theorem}{Theorem}[section]
\newtheorem{theo}[theorem]{Theorem}
\newtheorem{thm}[theorem]{Theorem}
\newtheorem{lemma}[theorem]{Lemma}
\newtheorem{lemm}[theorem]{Lemma}
\newtheorem{prop}[theorem]{Proposition}
\newtheorem{coro}[theorem]{Corollary}
\newtheorem{corollary}[theorem]{Corollary}
\newtheorem{fact}[theorem]{Fact}
\newtheorem{claim}[theorem]{Claim}
\theoremstyle{definition}
\newtheorem{defi}[theorem]{Definition}
\newtheorem{definition}[theorem]{Definition}
\newtheorem{rema}[theorem]{Remark}
\newtheorem{remark}[theorem]{Remark}
\newtheorem{exam}[theorem]{Example}
\newtheorem{example}[theorem]{Example}
\newenvironment{enonce*}[2][]{\par\medskip\noindent\textbf{#2.}\itshape}{\par\medskip}
\providecommand{\eg}{e.g.\ }
\providecommand{\ie}{i.e.\ }
\providecommand{\resp}{resp.\ }
\providecommand{\cf}{cf.\ }
\providecommand{\longthanks}[1]{\par\smallskip\noindent #1\par\smallskip}
\newtheorem{proposition}[theorem]{Proposition}
\newtheorem{conjecture}[theorem]{Conjecture}
\numberwithin{equation}{section}
\newcommand{\R}{{\mathbb R}}
\newcommand{\Z}{{\mathbb Z}}
\newcommand{\N}{{\mathbb N}}
\newcommand{\expref}[2]{\texorpdfstring{\hyperref[#2]{#1~\ref{#2}}}{#1~\ref{#2}}}
\newcommand{\epfmtdoi}[1]{\href{https://doi.org/#1}{doi:#1}}
\newcommand{\epfmt}[2]{#1:\,#2}
\newcommand{\cclass}[1]{\textsf{#1}}
\newcommand{\GF}{\operatorname{GF}}
\newcommand{\dproblem}[1]{\textsc{#1}}
\DeclareMathOperator{\E}{E}
\begin{document}
\begingroup\raggedright
\section*{1566. Deterministic approximation at the sharp satisfiable-system threshold for bounded-degree Boolean-field equations}
Johan Håstad. \emph{Satisfying Degree-d Equations over GF[2]n}. Theory of Computing volume 9 article 27 (2013), official complete journal PDF acquired 2026-09-30; canonical DOI 10.4086/toc.2013.v009a027 verified against official landing and printed PDF header. \url{https://theoryofcomputing.org/articles/v009a027/}. CC BY3.0.
\paragraph{Selected problem and scope (editorial).} Exact original Theorem4.1 only: a polynomial-time deterministic algorithm for m simultaneously satisfiable fixed-degree d Boolean-field equations returns an assignment satisfying at least ceil((2\textasciicircum{}\{1-d\}-2\textasciicircum{}\{1-2d\})m). Full fixed-degree dense-input and Boolean-quotient conventions, weight theorem statements, implied-affine-factor elimination, identity/linear-system recovery, complete AppendixA weight-bound induction with every factor-dependence case and final low-degree case, cited exact Viola generator statement and integer-gap rounding closing are retained. Gaussian elimination, affine-coordinate changes, elementary Boolean-polynomial operations and the named Viola polynomial generator are prerequisites; the source supplies the entire nontrivial local weight argument rather than omitting it behind a citation. Matching hardness and lifted-predicate outcomes are unselected and uncounted.
\paragraph{Difficulty (editorial).} After ordinary Boolean-polynomial and linear-algebra prerequisites, the deterministic threshold still depends on the complete source-supplied weight induction in AppendixA, including normalization of nonunique affine factors and all dependent/independent factor-family cases. The authored proof must also recover unknown implied affine factors efficiently from a polynomial identity, eliminate them while retaining satisfiability, and use the stated generator plus an exact integer gap to reach the rounded guarantee. These full structural-weight and algorithmic-closing obligations form one H2 task above elementary random-assignment or Gaussian-elimination qualification exercises.
\paragraph{Formatting (editorial).} Exact original human mathematical body and order retained; standalone typography, counter restoration and declared noncore printed locators only. Original comment prose excluded with percent guards retained; source typography and counters restored without upstream class execution. No mathematical repair or reconstruction.
\paragraph{Original source quirks (editorial).} The original Theorem2.2 display uses Q(X) while adjoining text uses x, and the complete appendix explicitly supplies its maximal independent affine-factor convention at1298–1303; both are retained. The appendix leaves the ordinary affine-hull/codimension characterization to the reader at1314–1322, with the full new weight induction nevertheless printed. Original wording including x\textasciicircum{}0/x\_0 at1406/1410, repeated sum-of wording and probability Pr typography is preserved without mathematical correction or new proof.
\endgroup
\clearpage
\setcounter{section}{0}
\section{Introduction}

The study of polynomial equations is a basic question
of mathematics.  In this paper we study a problem we
call \dproblem{Max-$d$-Eq}
where we are given a system of $m$ equations of degree
$d$ in $n$ variables over $\GF[2]$.  As we consider the
case of $d$ constant, all polynomials are given in
the dense representation, \ie, by a list of monomials.
Many problems can be coded as polynomial equations and in particular
it is easy to code \dproblem{3-Sat} as equations of degree 3
and thus determining whether we can simultaneously
satisfy all equations is \cclass{NP}-complete.   It is hence
natural to study
the question of satisfying the %
maximum
number of
equations and our interest turns to approximation
algorithms.  We say that an algorithm is a $C$-approximation
algorithm if it always returns a solution which satisfies
at least $C \cdot \mbox{OPT}$ equations where
$\mbox{OPT}$ is the number of equations satisfied
by the optimal solution.  
\setcounter{section}{1}
\section{Preliminaries}\label{sec: prelim}

We are interested in polynomials over $\GF[2]$.
Most polynomials we use are of
degree $d$ but also polynomials of degree one, that
we call ``affine forms'' play a special role.
We are not interested in the polynomials as formal
polynomials, but rather as functions mapping
$\GF[2]^n$ to $\GF[2]$ and hence we freely
use that $x_i^2=x_i$ and thus any
monomial in our polynomials can be taken to be multi-linear.
We start with the following standard result which we,
for completeness, even prove.

\begin{theorem}\label{thm basic deg d}
Any multivariate polynomial $P$ of degree
$d$ that is nonzero takes the
value 1 for at least a fraction
$2^{-d}$ of the inputs.
\end{theorem}


\begin{proof}
The proof is by induction over $n$ and $d$,
with the base case of $d=1$ which is true
as each linear polynomial is unbiased.

For the induction step, suppose $P(x)=P_0(x)+x_1 P_1(x)$
and let us consider what happens for the two possible
values of $x_1$.
If both $P_0$ and $P_0+P_1$ are non-zero we are
done by induction.  If not, as $P_1$ is of degree
at most $d-1$, the polynomial of $P_0$ and $P_0+P_1$ 
that is non-zero is of degree at most $d-1$.
Hence this polynomial takes the value 1 for at least a fraction $2^{1-d}$
of \emph{its} inputs.  As the set of inputs of this
polynomial constitutes half of the inputs of $P$,
the result follows also in this case.
\end{proof}

It is not difficult to see that this result is
tight by considering $P(x)=\prod_{i=1}^dx_i$,
or, more generally, products of $d$ linearly independent
affine forms.  
It is important for us
that these are the only cases of tightness.
This follows from a characterization by Kasami
and Tokura~\cite{kt70} of all polynomials
that are non-zero for at most a fraction
$2^{1-d}$ of the inputs.  
A consequence of their characterization is
the following theorem. 

\begin{theorem}\label{basic theorem}
Let $P$ be a %
degree-$d$
polynomial
over $\GF[2]$ which factors
as
$$P(x)= Q(X) \prod_{i=1}^r A_i (x) 
$$
where $A_i$ are linearly independent
affine forms and $Q$ does not
contain any affine factor.
Then the fraction of points on which $P(x)=1$
is at least
$$
2^{-r} (2^{1-(d-r)}-2^{1-2(d-r)})\,,$$
if $d\not =r$ and $2^{-d}$ if $d=r$.
\end{theorem}

Given that we do not need their full characterization
and the proof of the part that we need is shorter
than the proof of the full characterization, we
give the proof of \expref{Theorem}{basic theorem} in
an appendix.
\setcounter{section}{3}
\section{Completely satisfiable systems}\label{sec: p}

When studying systems where it is possible
to simultaneously satisfy all equations
the situation changes and let us start with
the positive results.

\subsection{Finding a good assignment}

Suppose we have an
equation of the form
$P(x)=1$ and that
this equation implies the affine
condition $A(x)=1$.  Then,
as the system is satisfiable,
we can use this equation to eliminate
one variable from the system, preserving the
degrees of all equations.  This is done
by taking some variable $x_i$ that appears
in $A$ and replacing it by $x_i+A(x)+1$ (note that
this function does not depend on $x_i$ as the two
occurrences of this variable cancel).
This substitution preserves the satisfiability of the
system and the process stops only 
when none of the current equations implies an affine condition.

Using \expref{Theorem}{basic theorem} we see
that when this process ends each equation
is satisfied by at least a fraction
$2^{1-d}-2^{1-2d}$ of the inputs and thus the following
theorem should not come as a surprise.

\begin{theorem}\label{thm upper sat}
There is a polynomial time algorithm that,
given a system of $m$ simultaneously satisfiable equations
of degree $d$ over $\GF[2]$, finds an
assignment that satisfies at least
$\lceil (2^{1-d}-2^{1-2d})m\rceil$ equations.
\end{theorem}

\begin{proof}
There are two points in the outlined argument
that require closer inspection.  The
first is the question of how to actually determine whether a 
polynomial equation implies an affine condition
and the second is to make sure that once
the process of finding implied affine conditions
has ended we can indeed deterministically find a
solution that satisfies the expected number of
equations.
Let us first address the issue of determining
whether a given equation implies an affine
condition.

Suppose $P(x)=1$ implies $A(x)=1$
for some unknown affine function $A$.
Let us assume that $x_1$ appears in $A$ with
a nonzero coefficient.  
We may write
$$
P(x)=P_0(x)+P_1(x)x_1$$
where neither $P_0$ nor $P_1$ depends on $x_1$.  We want
to prove that $P(x)=A(x)P_1(x)$ and towards this consider
\begin{eqnarray}\label{eq pq}
Q(x)=P(x)+A(x)P_1(x)\,.
\end{eqnarray}
As $x_1$ appears with coefficient one in $A$ it follows
that $Q$ does not depend on $x_1$ and let us assume that
$Q$ is not identically 0.  Choose any values
for $x_2, x_3 \ldots x_n$, to make $Q(x)=1$
and set $x_1$ to make $A(x)=0$.  It follows
from \eqref{eq pq} that $P(x)=1$ 
and thus we have found a counterexample 
to the assumed implication.  We can hence conclude
that $Q\equiv 0$ and we have 
$$
P(x)=A(x) P_1(x)\,.
$$

We claim furthermore that this procedure is entirely
efficient.  Namely given $P$ and the identity of
one variable occurring in $A$, $P_1$ is uniquely
defined.  Once $P_1$ is determined the rest
of the coefficients of $A$ can be found by
solving a linear system of equations.
As there are only $n$ candidates for a variable
in $A$ and solving a linear system of equations
can be done in polynomial time we conclude that the entire
process of finding implied affine conditions
can be done in polynomial time.

Once this process halts we need to find an assignment that satisfies the
expected number of equations.  This could
possibly be done in a manner similar
to how it was done in the proof of Theorem~3.1
but for the general case we do not know how to, in deterministic 
polynomial time,
find a suitable bound for the probability that an equation
is satisfied.  It is true that the bound of \expref{Theorem}{basic theorem}
can be calculated deterministically, but as we are not able to
prove it by induction, setting only one variable at the time,
it is not clear how to use it to guide an algorithm.
We turn to a different approach and we need the following
result by Viola~\cite{viola09}.

\begin{theorem}[\cite{viola09}]\label{thm:viola}
For any constant $d$ and $\epsilon > 0$ there exists
a polynomial time pseudo-random generator $G$ with
seed length $O(d \log n + d 2^d \log (1/\epsilon))$
and output in $\{ 0,1\}^n$ such that
$$
|\Pr[P(x)=0]- \Pr[P(G(y))=0]| \leq \epsilon
$$
for any polynomial $P$ of degree $d$.  Here $x$ is 
a uniformly random string in $\{ 0,1\}^n$ and
$y$ is a uniformly random seed.
\end{theorem}

This theorem implies that if we are willing
to sacrifice a small $\epsilon$ it is enough
to consider the candidate assignments 
$G(y)$ for all possible seeds $y$ and taking the assignment
from this set that satisfies the %
maximum
number of equations.  Let us be more precise.

We know that the expected number of satisfied
equations when we consider a uniformly random
input is $(2^{1-d}- 2^{1-2d}) m$ and, by
\expref{Theorem}{thm:viola}, the same number, when
we consider a random output from $G$, is
at least $(2^{1-d}- 2^{1-2d}-\epsilon) m$.
Setting $\epsilon= 2^{-2d} m^{-1}$ and observing that
the maximum number of equations satisfied is
an integer we see that this number is at least
$$
\lceil (2^{1-d}- 2^{1-2d}) m-2^{-2d} \rceil \,.
$$
Finally, as $(2^{1-d}- 2^{1-2d}) m$ is an integer
multiple of $2^{1-2d}$, we have that
$$
\lceil (2^{1-d}- 2^{1-2d}) m-2^{-2d} \rceil 
=\lceil (2^{1-d}- 2^{1-2d}) m \rceil\,,
$$
and the proof is complete.
\end{proof}

Note that the generator of Viola is only needed to
derandomize the algorithm and had we been content
with a randomized algorithms we could simply replace
the last step by random sampling, obtaining a much
more efficient algorithm.

\appendix
\section{Proof of \protect{\expref{Theorem}{basic theorem}}}

The goal of the current section is the prove 
\expref{Theorem}{basic theorem}.
We remind the reader that this result
follows from the characterization by
Kasami and Tokura~\cite{kt70} of 
all codewords of the Reed-Muller code
that have weight at most twice the
minimal weight.

First note that the bound of \expref{Theorem}{basic theorem} is
sharp as it is obtained by
$$x^\alpha (x^\beta+x^\gamma)$$ where $\alpha$,
$\beta$, and $\gamma$ are disjoint multi-indices
where the size of $\alpha$ is $r$ and the sizes
of $\beta$ and $\gamma$ both are $d-r$.

\begin{proof}
We prove the statement by induction and
we establish $d=2$ as the base case below
to avoid degenerate cases.  The following
easy observations are useful for us.

\begin{itemize}
\item As $d-r$ cannot equal 1, it follows
that for any degree-$d$ polynomial that
is not a product of affine factors, the
fraction of inputs on which it takes
the value one is at least $(3/2)\, 2^{-d}$.

\item The lower bound to be proved never exceeds 
$2^{1-d}$, and thus this bound is
always sufficient (but not always possible).
\end{itemize}
The statement for general $d$ and $r$ follows
from the case of $d-r$ and $0$ and hence we
may focus on the case of $r=0$ (but of course
use arbitrary $r$ in the inductive statements).

A fact that needs to be taken into account is, as we
are using that $x_i^2=x_i$,
that the factorization is not unique.  In
particular if $A$ and $A'$ are
two affine factors of a polynomial $P$ then another
factor is $1+A+A'$
as 
$$
(1+A+A')A=A'A\,.
$$
Thus if we have several affine factors
we can construct new affine factors by taking
the sum of of an even number of such factors
added with the constant 1 or the sum of
an odd number of such factors. 
Let us point out that the statement
of \expref{Theorem}{basic theorem} that $Q$ does not
contain any affine factor is intended to mean that there is no
factorization of $P$ that contains the given
affine factors and one more linearly independent
affine factor.

Another useful
fact is that an affine function $A(x)$ is a factor of 
a polynomial $P$
iff $A(x)=0$ implies $P(x)=0$ or equivalently
if $P(y)=1$ implies $A(y)=1$.  The non-obvious
direction of this statement was established during
the proof of \expref{Theorem}{thm upper sat}
when identifying implied affine constraints.

For the reader worried about the non-unique factorization,
let us point out the following fact that we leave to
the reader to verify.  The number of affine factors
is the co-dimension of the affine hull of all
points such that $P(x)=1$.  The factors that
may appear in the factorization are the affine equations
defining this affine hull.  In the full factorization
we may take any full set of linearly independent
equations.

Let us address the case of $d=2$, which can
be established by the normal form of 
%
degree-two 
polynomials, but let us follow a
different path to prepare for the general proof.
As stated above, the interesting case is
when $P$ has no affine factor and let
us write $P(x)=P_0(x)+x_1P_1(x)$.

Consider setting $x_1$ to its two values and as
$P$ does not contain an affine factor neither
of the induced polynomials can be identically 0.  Furthermore
if neither of these settings result in a polynomial
with an affine factor we are done by
induction.  Let us finally assume that $x_1=0$
results in an affine factor, which we, by an
affine change of coordinates, can assume is $x_2$.
Thus we can assume that 
$$
P(x)= x_2 A_2(x)+x_1A_1(x)\,,
$$
for two affine functions $A_1$ and $A_2$ where
we also can assume that $x_i$ does not appear in $A_i(x)$ as it
can be replaced by 1 giving the same result.
The main case is that the collection of $x_1,x_2, A_1(x),$
and $A_2(x)$ form linearly independent affine functions and
in this case the fraction of inputs for which
$P$ is one is exactly $3/8$ which is the claimed bound.
We have a number of cases to consider when the four
functions are not linearly independent.

\begin{enumerate}
\item $A_1(x) \equiv 1$.  If $A_2(x)=x_1$, then $x_1$ is a factor
of $P$ while if $A_2(x)=1+x_1$ then $P$ is one with
probability $3/4$. Finally if $A_2(x)$ is linearly independent
of $x_1$ this probability is $1/2$.  

\item $A_1(x)=x_2$ makes $x_2$ a factor of $P$.

\item $A_1(x)=1+x_2$ makes $Pr[P(x)=1]=1/2$ unless we
have $A_2(x) \equiv 1$ when this probability is $3/4$.

\item $A_1(x)$ is linearly independent of $x_2$.  In this case,
by an affine change of variables, we can assume that $A_1(x)=x_3$.
Now since $A_2(x)$ does not contain $x_2$, is linearly
dependent of $x_1$ and $x_3$, and is not a factor
of $x_1x_3$ (ruling out also $(1+x_1+x_3)$), 
$A_2(x)$ must equal one of the functions 1, $1+x_1$
or $1+x_3$.  In each of these cases it is easy
to check that $Pr[P(x)=1]$ is $1/2$.

\end{enumerate}
This finishes the case $d=2$ end we turn to the
general case.  Not surprisingly, also here we
end up analyzing a number of cases.

As in the case $d=2$ we can assume that $P$
has no affine factors but the polynomial resulting
when substituting $x_1=0$ gives a polynomial
with at least one affine factor.  Picking
a full set of linearly independent factors
and making an affine transformation we can
assume that
$$
P(x)=x_2 x^\beta P_2(x)+ x_1 P_1(x)$$
where $\beta$ is a  possible empty multi-index and
$P_2$ has no affine factors. First
let us consider affine factors in $P_1$.  

Let $\prod_{i=1}^r A_i(x)$ be the affine
factors that appear in $P_1$.  We have two
cases depending whether each $A_i$ that might appear
in the factorization, together with the
affine forms that appear in the first product
(i.\,e., the coordinate functions given by $x_2$
and the elements of $\beta$) are linearly independent.

Suppose these functions are not linearly independent and hence that
we can choose the factorization such
that $A_1(x)$ only depends
on $x_2$ and the variables in $\beta$.  Let
us look at the point $x^0$ where $x_2$ and all
elements of $\beta$ equals $1$.  If $A_1(x^0)=1$
then $A_1(x)$ is a factor of $P$
and this is a contradiction of the assumptions.
If, on the other hand $A_1(x_0)=0$ then the sets of points
where $x_2 x^\beta P_2(x)=1$ and $x_1 P_1(x)=1$
are disjoint and, as each of these sets is of density
at least $2^{-d}$,  we get $\Pr[P(x)=1] \geq 2^{1-d}$ and the
theorem follows also in this case.

Now assume that any affine form  appearing in the factorization
of $P_1$ is linearly independent of $x_2$ and the variables
in $\beta$.  Then, by an affine transformation,
we can assume that 
\begin{eqnarray}
P(x)=x_2 x^\beta P_2(x)+ x_1 x^\gamma P_1'(x)\,,
\end{eqnarray}
where $\beta$ and $\gamma$ are disjoint
multi-indices and no more affine factors
can be pulled out of $P_2$ or $P_1'$.

Let us analyze what happens for
the four possible simultaneous assignments
of values to $x_1$ and $x_2$.
When both are 0 we get a function that is
identically 0 which is not good for us
but in the other cases we get non-zero polynomials
of degree $d-1$ and we now analyze the 
structure of these polynomials.

When $x_1=0$ and $x_2=1$
we get $x^\beta P_2(x)$, when
$x_1=1$ and $x_2=0$ we get
$x^\gamma P_1'(x)$, and in the final
case we have
$$
W(x)=x^\beta P_2(x)+ x^\gamma P_1'(x)\,.$$
Note that $W$ is not identically 0
as $(x_1+x_2)$ then would have been
an affine factor of $P$ and furthermore that
$x^{\beta}P_2$, $x^{\gamma}P_1'$, $W$ all are of degree at
most $d-1$.

Suppose first that neither $P_2$ nor $P_1'$ is
the constant one.
Then, using the bound that a 
%
degree-$(d-1)$
polynomial that
is not a pure product of affine factors is one with
probability at least $(3/2)\,2^{1-d}$, we have
that $\Pr[P(x)=1]$ is at least
$$
\frac 14 \left(\frac 32 2^{1-d}+\frac 32 2^{1-d}+2^{1-d}\right)= 2^{1-d}
$$
proving the bound in this case.
By symmetry we may thus assume 
that $P_2\equiv 1$.
If $\gamma=\emptyset$
or $W$ does not contain any affine factor,
then $\Pr[P(x)=1]$ is at least
$$
\frac 14 (2^{1-d}+2^{1-d}+2^{2-d}-2^{3-2d})\,,
$$
which exactly equals the claimed bound.
Thus we can assume that $\gamma$ is non-empty
and $W$ contains an affine factor $A(x)$ and remember
that \begin{eqnarray}\label{weq}
W(x)=x^\beta + x^\gamma P_1'(x)\,.
\end{eqnarray}

Suppose $A$ only depends
on variables in $\beta$.  If it is fixed to
1 by setting all variables in $\beta$
to one, then $A(x)$ is a factor of $x^\beta$
and hence also of $W(x)+x^\beta=x^\gamma P_1'(x)$
and hence also of $P(x)$, contradicting assumptions.
If $A(x)$ is forced to 0 by this assignment
then setting any variable
in $\gamma$ (remember it is non-empty and disjoint
from $\beta$) to 0 and we get $W(x)=0$ while
$x^\beta=1$ and $x^\gamma P_1'(x)=0$ contradicting
\eqref{weq}.  Thus we can assume
that $A(x)$ depends on some variable
outside $\beta$.

If we can fix some variable in $\gamma$ to 0,
the variables of $\beta$ to one and $A$ to zero
we get a contradiction to \eqref{weq} as
$x^\beta=1$ while $W(x)=0$ and $x^\gamma=0$.
If the size of $\gamma$ is
at least 2 or $W$ contains at least two different
affine factors we claim that this must be possible.  Suppose first
that the size of $\gamma$ is at least 2.

As $A(x)$ does not only depend on variables
in $\beta$ we can fix these variables to
one, and then pick a suitable variable
in $\gamma$ to fix to 0 without fixing
the value of $A(x)$. We can then fix
additional variables to make $A(x)=0$
obtaining the desired contradiction.

Now suppose that $W$ contains at least
2 affine factors and let $A'$ be a factor distinct
from $A$.  In this case,
one of $A$, $A'$ and $1+A+A'$ is
a factor of $W$ and does not
depend on the first variable of $\gamma$
(and some variable outside $\beta$).
It follows again that we can
make $x^\beta=1$, $x^\gamma=0$ and $W(x)=0$,
again obtaining a contradiction.

The only remaining case is when $\gamma$ is
of size one and $W$ has one affine factor.
In this case, by induction
$\Pr[P(x)=1]$ is at least
$$
\frac 14 \left(2^{1-d}+ 2\cdot \frac 12 \left(2^{3-d}-2^{5-2d}\right)\right).
$$
For $d$ at least 4 this is at least $2^{1-d}$
and we are done.  Finally note that in the
case $d=3$, $P_1'$ as well as the co-factor
of $A$ in $W$ are of degree at most one
and hence if they do not contain an
additional factor they must be the constant
1 and the theorem follows also in this case.
\end{proof}
\clearpage
\begingroup\raggedright
%
%
%
%
%
\providecommand{\bibhead}[1]{}
%
\expandafter\ifx\csname pdfbookmark\endcsname\relax%
  \providecommand{\tocrefpdfbookmark}{}
\else\providecommand{\tocrefpdfbookmark}{%
   \hypertarget{tocreferences}{}%
   \pdfbookmark[1]{References}{tocreferences}}%
\fi

\tocrefpdfbookmark
\begin{thebibliography}{10}

\bibitem{almss}\bibhead{almss}
{\sc Sanjeev Arora, Carsten Lund, Rajeev Motwani, Madhu Sudan, and Mario
  Szegedy}: Proof verification and the hardness of approximation problems.
\newblock {\em J. ACM}, 45(3):501--555, 1998.
\newblock Preliminary version in
  \href{http://dx.doi.org/10.1109/SFCS.1992.267823}{FOCS'92}. See also at
  \href{http://eccc.hpi-web.de/report/1998/008/}{ECCC}.
\newblock [\epfmtdoi{10.1145/278298.278306}]

\bibitem{ah12}\bibhead{ah12}
{\sc Per Austrin and Johan H{\aa}stad}: On the usefulness of predicates.
\newblock {\em ACM Trans. Computation Theory}, 5(1):1:1--1:24, 2013.
\newblock Preliminary version in
  \href{http://dx.doi.org/10.1109/CCC.2012.18}{CCC'12}.
\newblock [\epfmtdoi{10.1145/2462896.2462897}]

\bibitem{AustrinMossel:08}\bibhead{AustrinMossel:08}
{\sc Per Austrin and Elchanan Mossel}: Approximation resistant predicates from
  pairwise independence.
\newblock {\em Comput. Complexity}, 18(2):249--271, 2009.
\newblock Preliminary version in
  \href{http://dx.doi.org/10.1109/CCC.2008.20}{CCC'08}. See also at
  \href{http://eccc.hpi-web.de/report/2008/009/}{ECCC}.
\newblock [\epfmtdoi{10.1007/s00037-009-0272-6}]

\bibitem{bgs98}\bibhead{bgs98}
{\sc Mihir Bellare, Oded Goldreich, and Madhu Sudan}: Free bits,
  \protect{PCPs}, and nonapproximability---towards tight results.
\newblock {\em SIAM J. Comput.}, 27(3):804--915, 1998.
\newblock Preliminary version in
  \href{http://dx.doi.org/10.1109/SFCS.1995.492573}{FOCS'95}. See also at
  \href{http://eccc.hpi-web.de/report/1995/024/}{ECCC}.
\newblock [\epfmtdoi{10.1137/S0097539796302531}]

\bibitem{jhacm}\bibhead{jhacm}
{\sc Johan H{\aa}stad}: Some optimal inapproximability results.
\newblock {\em J. ACM}, 48(4):798--859, 2001.
\newblock Preliminary version in
  \href{http://dx.doi.org/10.1145/258533.258536}{STOC'97}.
\newblock [\epfmtdoi{10.1145/502090.502098}]

\bibitem{jhdegd}\bibhead{jhdegd}
{\sc Johan H{\aa}stad}: Satisfying degree-$d$ equations over
  \protect{$GF[2]^n$}.
\newblock In {\em Proc. 14th Internat. Workshop on Approximation Algorithms for
  Combinatorial Optimization Problems (APPROX'11)}, pp. 242--253. Springer,
  2011.
\newblock [\epfmtdoi{10.1007/978-3-642-22935-0\_21}]

\bibitem{hps93}\bibhead{hps93}
{\sc Johan H{\aa}stad, Steven Phillips, and Shmuel Safra}: A well-characterized
  approximation problem.
\newblock {\em Information Processing Letters}, 47(6):301--305, 1993.
\newblock Preliminary version in
  \href{http://sceas.csd.auth.gr/php/conferences.php4?conf_id=6365}{ISTCS'93}.
\newblock [\epfmtdoi{10.1016/0020-0190(93)90076-L}]

\bibitem{kt70}\bibhead{kt70}
{\sc Tadao Kasami and Nobuki Tokura}: On the weight structure of
  \protect{Reed-Muller} codes.
\newblock {\em IEEE Trans. Inform. Theory}, 16(6):752--759, 1970.
\newblock [\epfmtdoi{10.1109/TIT.1970.1054545}]

\bibitem{razmosh10}\bibhead{razmosh10}
{\sc Dana Moshkovitz and Ran Raz}: Two-query {PCP} with subconstant error.
\newblock {\em J. ACM}, 57(5):29:1--29:29, 2010.
\newblock Preliminary version in
  \href{http://dx.doi.org/10.1109/FOCS.2008.60}{FOCS'08}. See also at
  \href{http://eccc.hpi-web.de/report/2008/071/}{ECCC}.
\newblock [\epfmtdoi{10.1145/1754399.1754402}]

\bibitem{r}\bibhead{r}
{\sc Ran Raz}: A parallel repetition theorem.
\newblock {\em SIAM J. Comput.}, 27(3):763--803, 1998.
\newblock Preliminary version in
  \href{http://dx.doi.org/10.1145/225058.225181}{STOC'95}.
\newblock [\epfmtdoi{10.1137/S0097539795280895}]

\bibitem{viola09}\bibhead{viola09}
{\sc Emanuele Viola}: The sum of {$D$} small-bias generators fools polynomials
  of degree {$D$}.
\newblock {\em Comput. Complexity}, 18(2):209--217, 2009.
\newblock Preliminary version in
  \href{http://dx.doi.org/10.1109/CCC.2008.16}{CCC'08}.
\newblock [\epfmtdoi{10.1007/s00037-009-0273-5}]

\end{thebibliography}
\endgroup
\end{document}
